% !TeX root = ../papers/nc-p3.tex
\section{二直線列のない純関係自由な分岐の位置と長さ}
\label{sec:branch-centers-opposite-phases}
\begin{theorem}\label{thm:branch-centers-opposite-phases}
曲面射影と純関係のない標準正則窓で、片側分岐があり、二直線列がないとする。
特に $9\le\rho\le11$ の残余分岐はこの範囲にある。
四位置のうち相対する二位置だけに順中心が一つずつ、同じ二位置に逆中心が一つずつある。
二つの曲線列の動く位置数の組は順序を除いて
\[
                           (3,3),\quad(3,7),\quad(5,5)
\]
に限る。これは必要条件であり、この三配置の存在を主張しない。
\end{theorem}
\begin{proof}[位置の制約と次数十の接空間]
位置を法4で読む。$F_i,B_i$ を位置 $i$ の順・逆中心集合とする。
$F_i$ の一つの中心 $p$ の後続直線を $L_p$、その反対三次曲線を $C_p$ とする。
命題\ref{prop:pure-free-branch-cubic}より辺 $i+1\to i+2$ の動く二成分間の移動は、
$L_p$ からのもの、または $C_p$ に着くものに限る。
一般前点・後点の一意性によりこの相異なる移動は高々二本である。
特に $F_i$ が二点なら両移動は二つの初回の移動で既に尽きる。
二中心の後続直線は交わらない。どの点でも前点は元の中心一つだからである。
$B_{i-1}$ が存在すれば、その直前の動く二成分の移動はこの同じ辺で第二成分が直線なので、
第一成分は $L_p$ に等しい。その一般前点は $p$ で、そこから二直線列になってしまう。
従って $F_i$ があれば $B_{i-1}$ は空である。また $B_{i+2}$ の任意の中心の反対三次曲線に
$F_i$ の全中心が入る。逆向きの二者択一でこれに入らない順中心の後続直線は、
当該逆中心の前点直線に等しくなり、動く位置が一つだけの列を与えてしまうからである。

ある位置で相異なる動く曲線の次数和が10に達するとき、そこにある順中心を通る
それらの曲線の解析的な枝が二本以下で全て滑らかであることは不可能である。
実際、順辺の一般五式の完全交叉は次元1で、その射影は五行四列評価行列の最大小行列式で切られる。
次数を数える曲線を含み、台の次元1なので高さ2のHilbert--Burch解消
$0\to4\mathcal O(-5)\to5\mathcal O(-4)\to\mathcal I\to0$ を持ち、
埋込点のない次数10の曲線である。次数を使い切った各成分は被約で、全スキームがその被約和となる。
一方、順中心では評価階数2なので四次小行列式の一次微分は全て零で、接次元は3である。
滑らかな一枝の接次元は1、滑らかな二枝の和も高々2である。
後者の完備局所環は二つの $k[[t]]$ を共通の交叉商 $k[[t]]/(t^e)$ 上で貼った環で、
パラメータを合わせれば $(t,t),(0,t^e)$ で生成できる。
正規化で不分岐な射影直線の枝は完備局所環に一様化変数を含むので滑らかである。

順中心の位置が一つ $i$ だけと仮定する。各列で内部に順中心へ戻る必要があるため
動く位置は少なくとも五つである。最初の三次数は次数公式で $1,3,6$ となる
（次の二位置の順側欠損が零、$b_1=0$）。中心二つなら相異なる二曲線が同じ位置で
ともに次数6となり上限10に反する。一つなら第四次数は $10+b_2\ge10$ で位置 $i$ の上限を使い切る。
列に同じ成分は再登場しないので内部で順中心を訪れるのはこの成分だけである。
命題\ref{prop:terminal-branch-contact}からそのパラメータは一つで不分岐であり、
先の接空間に反する。逆向きにも占有位置が一つだけの場合はない。
$F_i$ と $B_{i-1}$ の排他性より、順・逆とも正確に二位置ずつを占める。

\end{proof}
\begin{proof}[隣り合う位置の排除]
$F_0,F_1$ が存在すれば $B_1,B_2$ が存在する。$F_0$ は一つである（二つなら $F_1$ は空）。
$F_1$ の全中心は $L_{p_0}$ 上にある。二つならその上の $f_1\ge2$ から第二次数は1以下で、
二直線列になってしまう。逆側も一つずつである。二列の最初の次数は
$p_0$ の側で $1,2,3$、$p_1$ の側で $1,3,6$ である。
前者の第二次数は正の $f_1$ と二直線の排除による。第三の像は $F_1$ の二者択一により三次曲線である
（第二の円錐曲線は $F_1$ の後続直線でない）。第三像が定点なら終中心が位置3に来てしまう。
後者は次の二位置に順中心がないことによる。いずれも少なくとも四つの動く位置がある。
$p_0$ の第四次数は $4+b_2$、同位置0の $p_1$ の第三次数は6なので $b_2=0$ で次数10を使い切る。
$p_1$ の第三の曲線上に順中心 $p_0$ がなければ、第四次数は $10+b_2\ge10$ となり、
位置1にある $p_0$ の最初の直線と両立しない。
従って両中心は互いに到達し、唯一の終端強連結成分をなす。
位置0で $p_0$ を通る二成分の各々は、多くとも一つのパラメータでそこに至り不分岐である。
これも次数10の接空間に反する。

\end{proof}
\begin{proof}[相対する二位置]
添字を命題の表示のために選び、順・逆とも位置0と2を占めるとする。
列長は奇数で3以上である。$F_0$ が二つなら辺 $1\to2$ の移動枠が既に埋まるので、
$F_2$ の全列は長さ3、$F_0$ の列は長さ3または5である。
長さ5のものがなければ $B_0$ は二つで、$F_0$ の列の中央にある二本の三次曲線は
逆向きの二中心の二者択一からちょうど $B_0$ の反対三次曲線二本である。
上で示したことから $F_2$ の中心は両方に属する。しかしその点の前点は一意であり、
交わらない $F_0$ の二つの後続直線の両方には入れない。
長さ5のものがあれば $B_0$ が空でないため残る一つは長さ3で、
$B_2$ の上限二つにより $F_2$ は一つである。二つの $B_2$ の反対三次曲線の一つは
$F_2$ の長さ3の列の中央曲線であり、そこには $F_0$ の二点がともに属する。
一方、$F_0$ の長さ3の列では初めの $1,3,1$ から $f_2=5>0$。
長さ5の列でも最後から第二の次数が3である（逆中心の隣の位置に逆中心がない）。
第三の次数を $d$ とすれば第四次数 $3=3d-8+b_2$ より $d\le3$、
ゆえに $f_2=6-d>0$。両方の $F_0$ が $F_2$ に戻り、全三中心が終端強連結成分となるが、
$F_2$ が相異なる二点へ二パラメータで進むのは不可能である。
これで二重の順中心を排除し、逆側も同じである。

最後に各枝の長さは9以下である。自分の最初の移動位相を再訪して動く曲線を渡れるのは
残る一枠だけであり、三回目には進めない。さらに両方の列が長さ5以上なら、
他の列も第一の移動位相を一回使うので、自身の長さは5を超せない
（奇数長7以上なら自分の第五曲線から第六曲線へ渡る）。
両終端は別々の位置を占めるので二つの長さはともに $3\bmod4$ またはともに $1\bmod4$ となる。
これらの条件から表示の三組に限られる。
\end{proof}
